\scanpage{7}
\uline{Unités} On définit

\[
\deleted{\unsure{\delta}}\Delta_X\in\operatorname{Hom}_{\unsure{V|\mathcal V}^{\mathcal A}\unsure{{}_{c}|{}}}(X,X)=A(X\times X)
\]

en considérant

\[
\delta_X:X\mathbin{-}X\times X\qquad(\text{diag})
\]

et définissant

\[
\Delta_X=\deleted{\unsure{\delta_{X*}A(\delta_X)}}A(\delta_X)_*(1_{A(X)})
\]

Prouvons que $\Delta_X$ est une unité bilatère pour la composition : on a

\diagram{\begin{tikzcd}[ampersand replacement=\&]
X\arrow[r,dash,"\Delta_X"]\&X\arrow[r,dash,"u"]\&Y
\end{tikzcd}}

\[
u\circ\Delta_X=p_{31*}\bigl(p_{32}^*(u)\,p_{21}^*(\Delta_X)\bigr)
\]

\diagram{\begin{tikzcd}[ampersand replacement=\&,column sep=small]
\&X\times X\times Y\arrow[dl,"p_{21}"']\arrow[d,dash]\arrow[dr,"p_{32}"]\&\\
X\times X\&X\times Y\&X\times Y
\end{tikzcd}}

\[
p_{21}^*(\Delta_X)=p_{21}^{\unsure{'*|*}}(\delta_{X*}(1))\overset{\boxed{\text{échange}}}{=}q_*(p^*(1))=q_*(1)
\]

\[
\begin{aligned}
\text{donc}\quad p_{32}^*(u)\,p_{21}^*(\Delta_X)&=p_{32}^*(\unsure{u_X|u})\,q_*(1)
\underset{\substack{\text{form. de proj.}\\\text{Attention \uline{\unsure{signe|signes}} !}}}{=}q_*\bigl(q^*(p_{32}^*(u_X))\bigr)\\
\deleted{\unsure{\text{or}}}\quad&=q_*(u)\qquad\text{car }p_{32}q=\mathrm{id}_{X\times Y}
\end{aligned}
\]

\diagram{\begin{tikzcd}[ampersand replacement=\&]
X\times Y\arrow[r,"\delta_X\times\mathrm{id}_Y=q"]\arrow[d,"{\unsure{p|p_*}}"']\&X\times X\times Y\arrow[d,"p_{21}"]\\
X\arrow[r,"\delta_X"']\&X\times X
\end{tikzcd}}

\margin{\fbox{\parbox{2cm}{Signes $(-1)^{ux}$ ; pas de signes pour \unsure{constater|calculs} que $\Delta_X\circ v=v$ $v:Y\mathbin{-}X$}}}

\[
\begin{aligned}
\text{donc}\quad u\circ\Delta_X&=p_{31*}\,q_*(u)=u\\
&\text{car }\deleted{q}\,p_{31}q=\mathrm{id}
\end{aligned}
\]

On prouve de façon analogue que $\Delta_X\circ u=u$.

\uline{Donc on trouve bien une catégorie $V^{\mathcal A}$}

[ \uline{NB} Pour définir la composition on n'a eu besoin que des $p_*,p^*$ pour projections $X\times Y\underset{\deleted{\unsure{\leftarrow|\rightarrow}}}{\longrightarrow}X$ ; mais pour \deleted{définir} construire les unités on a utilisé aussi $q_*,q^*$ pour les imm. diagonales et leurs produits par un $Y$ \ldots{} ]

% archive sheet number: 6
